\begin{proposition}[4.7.8]
\label{IV.2.4.7.8-fr}
Bajo las hipótesis de \textup{(4.7.5)}, sea $x$ un punto maximal de
$\operatorname{Supp}(\mathscr F)$. Sea $k'$ una extensión perfecta de $k$~;
pongamos $X'=X\otimes_k k'$, $\mathscr F'=\mathscr F\otimes_k k'$~; entonces la
longitud geométrica $\lambda_x(\mathscr F)$ de $\mathscr F$ en $x$ es igual a la
longitud del $\mathcal O_{x'}$-módulo $\mathscr F'_{x'}$, en todo punto maximal
$x'$ de $\operatorname{Supp}(\mathscr F')$ por encima de $x$. Además, existe una
extensión radicial finita $k_1$ de $k$ tal que, poniendo
$X_1=X\otimes_k k_1$, $\mathscr F_1=\mathscr F\otimes_k k_1$,
$\lambda_x(\mathscr F)$ es igual a la longitud del
$\mathcal O_{x_1}$-módulo $(\mathscr F_1)_{x_1}$ en todo punto maximal $x_1$ de
$\operatorname{Supp}(\mathscr F_1)$ por encima de $x$.
\end{proposition}

\begin{proof}
La observación de (4.7.7) muestra que es posible limitarse al caso en que $x$ es
un punto maximal de $X$. Además, la cuestión es local sobre $X$, por lo que
se puede suponer $X=\operatorname{Spec}(A)$ afín noetheriano,
$\mathscr F=\widetilde M$, donde $M$ es un $A$-módulo de tipo finito~; sean
$B=A\otimes_k k'$, $\mathfrak p$ (resp. $\mathfrak q$) el ideal primo minimal
de $A$ (resp. $B$) correspondiente a $x$ (resp. $x'$). Se tiene entonces,
aplicando la fórmula (4.7.2.1),
\[
\operatorname{long}_{\mathcal O_{x'}}(\mathscr F'_{x'})
=\operatorname{long}_{\mathcal O_x}(\mathscr F_x)
 .\operatorname{long}_{B_{\mathfrak q}}
 (B_{\mathfrak q}/\mathfrak pB_{\mathfrak q}).
\]
Pero, como $\boldsymbol{k}(x)=A_{\mathfrak p}/\mathfrak pA_{\mathfrak p}$, el mismo
razonamiento que en (4.7.3) muestra que
$\operatorname{long}_{B_{\mathfrak q}}
(B_{\mathfrak q}/\mathfrak pB_{\mathfrak q})$ es igual a la longitud del anillo
local $(B/\mathfrak pB)_{\mathfrak q'}$, donde
$\mathfrak q'=\mathfrak q/\mathfrak pB$. El anillo
$(B/\mathfrak pB)_{\mathfrak q'}$ es un anillo local de
$(A/\mathfrak p)\otimes_k k'$ correspondiente a un ideal primo minimal de este anillo, y es,
por tanto, también un anillo local del anillo total de fracciones de
$\boldsymbol{k}(x)\otimes_k k'$; la longitud de tal anillo es, por definición,
$\mu_i(\boldsymbol{k}(x)\mid k)$, de donde resulta la primera afirmación. La segunda se
demuestra del mismo modo, reemplazando $k'$ por $k_1$ y utilizando (4.7.3, (iii)).
\end{proof}

Se dice también, en el caso en que $x$ es un punto maximal de
$\operatorname{Supp}(\mathscr F)$, que la longitud geométrica de $\mathscr F$ en
$x$ es la \emph{multiplicidad radicial del ciclo primo maximal} $\overline{\{x\}}$
asociado a $\mathscr F$.

\begin{corollary}[4.7.9]
\label{IV.2.4.7.9-fr}
Bajo las hipótesis de \textup{(4.7.5)}, sean $k'$ una extensión de $k$,
$X'=X\otimes_k k'$, $\mathscr F'=\mathscr F\otimes_k k'$~; sean $Z$ una
componente irreducible de $\operatorname{Supp}(\mathscr F)$ y $Z'$ una
componente irreducible de $\operatorname{Supp}(\mathscr F')$ que domina a $Z$
\textup{(4.2.7)}. Entonces las multiplicidades radiciales de $Z$ con respecto a
$\mathscr F$ y de $Z'$ con respecto a $\mathscr F'$ son iguales.
\end{corollary}

\oldpage[IV-2]{79}

\begin{proof}
Basta considerar una extensión perfecta $k''$ de $k'$ y, poniendo
$X''=X\otimes_k k''$, $\mathscr F''=\mathscr F\otimes_k k''$, un punto maximal
$x''$ de $\operatorname{Supp}(\mathscr F'')$ que está por encima del punto
genérico de $Z'$~; entonces se deduce de (4.7.8) que
$\operatorname{long}(\mathscr F''_{x''})$ es igual a las multiplicidades de $Z$
y de $Z'$.
\end{proof}

\begin{proposition}[4.7.10]
\label{IV.2.4.7.10-fr}
Sean $k$ un cuerpo, $X$ un $k$-preesquema localmente de tipo finito sobre $k$,
$\mathscr F$ un $\mathcal O_X$-módulo coherente, $x$ un punto de $X$ y $z_i$ los
puntos genéricos de las componentes irreducibles de
$\operatorname{Supp}(\mathscr F)$ que contienen a $x$. Para que $\mathscr F$ sea
geométricamente reducido en el punto $\boldsymbol{x}$ \textup{(4.6.22)}, es necesario y suficiente
que $x$ no pertenezca a ningún ciclo primo asociado inmerso de $\mathscr F$ y
que, para todo $i$, la longitud geométrica $\lambda_{z_i}(\mathscr F)$ de
$\mathscr F$ en el punto $z_i$ sea igual a $1$.
\end{proposition}

\begin{proof}
Demostremos primero que las condiciones del enunciado son suficientes. En
efecto, si $k'$ es una extensión finita de $k$, $x'$ un punto de
$X'=X\otimes_k k'$ por encima de $x$, y
$\mathscr F'=\mathscr F\otimes_k k'$, $x'$ no pertenece a ningún ciclo primo
asociado inmerso de $\mathscr F'$ en virtud de (4.2.7)~; además, en virtud de
(4.7.8), la longitud geométrica de $\mathscr F'$ es igual a $1$ en todo punto
genérico $z'_j$ de una componente irreducible de
$\operatorname{Supp}(\mathscr F')$ que contiene a $x'$, pues tal punto está por
encima de uno de los $z_i$ en virtud de (2.3.4)~; \emph{a fortiori}, se tiene
$\operatorname{long}(\mathscr F'_{z'_j})=1$, lo que demuestra que $\mathscr F'$
es reducido en el punto $x'$.

Recíprocamente, se deduce primero de (4.2.7) que, si para \emph{una} extensión finita
$k'$ de $k$, $\mathscr F'$ es reducido en todos los puntos $x'$ de $X'$ por
encima de $x$, entonces $x$ no pertenece a ningún ciclo primo asociado inmerso
de $\mathscr F$. Además, tomando como $k'$ el cuerpo $k_1$ del enunciado de
(4.7.8) (donde $x$ se reemplaza por uno de los $z_i$), la hipótesis de que
$\mathscr F$ es geométricamente reducido en el punto $x$ implica que
$\lambda_{z_i}(\mathscr F)=1$.
\end{proof}
